\documentclass[aspectratio=169,11pt]{beamer}

% Talk slides for: On Self-Similar Viscosity Models for Carreau Boundary Layer Flows
% Put this file in the same directory as the figures used by main.tex.
% Figure filenames are intentionally kept identical to the manuscript source.

\usetheme{Madrid}
\usecolortheme{seahorse}
\setbeamertemplate{navigation symbols}{}
\setbeamertemplate{caption}[numbered]

\usepackage{amsmath,mathtools,bm,booktabs}
\usepackage{graphicx}
\usepackage{appendixnumberbeamer}

\newcommand{\gd}{{\dot{\bm{\gamma}}}}
\newcommand{\Reyn}{\operatorname{Re}}
\newcommand{\De}{\operatorname{De}}
\newcommand{\dd}{\mathrm{d}}

% Compile-friendly figure macro: includes the manuscript figure if present;
% otherwise leaves a labeled placeholder so the deck still compiles.
\newcommand{\figinclude}[2][]{%
  \IfFileExists{#2}{\includegraphics[#1]{#2}}{%
    \fbox{\begin{minipage}[c][0.38\textheight][c]{0.82\linewidth}
    \centering\small Missing figure file:\\[0.4em]\texttt{\detokenize{#2}}
    \end{minipage}}%
  }%
}

\title[Self-Similar Models for Carreau Boundary Layers]{On Self-Similar Viscosity Models for Carreau Boundary Layer Flows}
\author[Reinberger, Barlow, Weinstein]{W. Cade Reinberger \and Nathaniel S. Barlow \and Steven J. Weinstein}
\institute[RIT]{Rochester Institute of Technology}
\date{}

\begin{document}

% 1
\begin{frame}
  \titlepage
\end{frame}

% 2
\begin{frame}{One-slide version}
\begin{columns}[T,onlytextwidth]
\begin{column}{0.58\textwidth}
\begin{itemize}
  \item Boundary-layer flows have high shear near the wall/leading edge and low shear far from the wall/downstream.
  \item Power-law and Newtonian models admit similarity solutions; Carreau rheology does not.
  \item Main question: when do simple self-similar models accurately predict Carreau wall shear?
  \item Main result: use power-law upstream, Newtonian downstream, with a computable transition criterion and quantified high-shear error.
\end{itemize}
\end{column}
\begin{column}{0.38\textwidth}
\centering
\figinclude[width=\linewidth]{rheo_thin_cropped_rotated.pdf}\vspace{0.4em}
\figinclude[width=\linewidth]{rheo_thick_cropped_rotated.pdf}
\end{column}
\end{columns}
\end{frame}

% 3
\begin{frame}{Roadmap for a 20-minute talk}
\begin{enumerate}
  \item Boundary-layer setup and viscosity models.
  \item Similarity reductions: power-law/Newtonian versus Carreau.
  \item New Sakiadis power-law solutions outside the previously accepted range.
  \item Wall-shear comparison and the transition parameter $\delta^*$.
  \item Why the high-shear power-law approximation has a finite offset.
\end{enumerate}
\end{frame}

% 4
\begin{frame}{Why the power-law model is both useful and incomplete}
\begin{columns}[T,onlytextwidth]
\begin{column}{0.48\textwidth}
\begin{block}{Power-law rheology}
\[
  \mu = K |\gd|^{\alpha-1}
\]
\begin{itemize}
  \item $\alpha<1$: shear-thinning.
  \item $\alpha>1$: shear-thickening.
  \item Analytically convenient; often self-similar.
\end{itemize}
\end{block}
\end{column}
\begin{column}{0.48\textwidth}
\begin{alertblock}{Problem}
\begin{itemize}
  \item It has no low-shear Newtonian plateau.
  \item In a boundary layer, low-shear and high-shear regions coexist in the same flow.
  \item So the far-field velocity profile can be wrong even when wall shear looks close.
\end{itemize}
\end{alertblock}
\end{column}
\end{columns}
\end{frame}

% 5
\begin{frame}{Carreau rheology: power-law behavior with Newtonian regularization}
\begin{columns}[T,onlytextwidth]
\begin{column}{0.48\textwidth}
\[
  \mu = \mu_0\left(1+(\lambda |\gd|)^2\right)^{(\alpha-1)/2}
\]
\vspace{-0.5em}
\[
  \mu \sim \mu_0 \quad (|\gd|\to 0),
  \qquad
  \mu \sim \mu_0\lambda^{\alpha-1}|\gd|^{\alpha-1}\quad(|\gd|\to\infty)
\]
\[
  |\gd|_c=\frac{1}{\lambda},\qquad K=\mu_0\lambda^{\alpha-1}.
\]
\end{column}
\begin{column}{0.48\textwidth}
\centering
\figinclude[width=0.94\linewidth]{rheo_thin_cropped_rotated.pdf}\vspace{0.2em}
\figinclude[width=0.94\linewidth]{rheo_thick_cropped_rotated.pdf}
\end{column}
\end{columns}
\end{frame}

% 6
\begin{frame}{Two canonical boundary layers}
\begin{columns}[T,onlytextwidth]
\begin{column}{0.52\textwidth}
\begin{itemize}
  \item \textbf{Blasius:} stationary wall, moving fluid.
  \item \textbf{Sakiadis:} moving wall, quiescent fluid.
  \item They are not equivalent by a frame shift because the governing boundary-layer equation is nonlinear.
\end{itemize}
\vspace{0.5em}
\begin{center}
\begin{tabular}{lll}
\toprule
Location & Blasius & Sakiadis \\
\midrule
$y=0$ & $u_x=0$ & $u_x=U$ \\
$y\to\infty$ & $u_x\to U$ & $u_x\to 0$ \\
\bottomrule
\end{tabular}
\end{center}
\end{column}
\begin{column}{0.44\textwidth}
\centering
\figinclude[width=\linewidth]{schematics_pdf_cropped.pdf}
\end{column}
\end{columns}
\end{frame}

% 7
\begin{frame}{Boundary-layer equations and comparison metric}
\begin{columns}[T,onlytextwidth]
\begin{column}{0.52\textwidth}
\begin{align*}
  &\frac{\partial u_x}{\partial x}+\frac{\partial u_y}{\partial y}=0,\\
  &\rho\left(u_x\frac{\partial u_x}{\partial x}+u_y\frac{\partial u_x}{\partial y}\right)
    =\frac{\partial \tau_{yx}}{\partial y},\\
  &\tau_{yx}=\mu(|\gd|)\frac{\partial u_x}{\partial y},
  \qquad |\gd|=\left|\frac{\partial u_x}{\partial y}\right|.
\end{align*}
\end{column}
\begin{column}{0.43\textwidth}
\begin{block}{Metric}
Wall shear / wall shear-rate is the central quantity:
\[
  \tau_i \quad (i=n,p,c).
\]
It is sensitive to the full boundary-value problem, but is still the most engineering-facing scalar output.
\end{block}
\end{column}
\end{columns}
\end{frame}

% 8
\begin{frame}{Power-law and Newtonian models are self-similar}
\begin{columns}[T,onlytextwidth]
\begin{column}{0.52\textwidth}
\[
  \eta = y\left(\frac{\rho U^{2-\alpha}}{Kx}\right)^{1/(\alpha+1)},
  \qquad u_x=Uf_p'(\eta;\alpha).
\]
\[
  \alpha(\alpha+1)f_p''' + f_p f_p''|f_p''|^{1-\alpha}=0.
\]
\end{column}
\begin{column}{0.43\textwidth}
\begin{block}{Boundary conditions}
\textbf{Blasius:}
\[
 f_p(0)=0,\quad f_p'(0)=0,\quad f_p'(\infty)=1.
\]
\textbf{Sakiadis:}
\[
 f_p(0)=0,\quad f_p'(0)=1,\quad f_p'(\infty)=0.
\]
\end{block}
\end{column}
\end{columns}
\end{frame}

% 9
\begin{frame}{Carreau is not self-similar, but reduces to a one-parameter family}
\begin{columns}[T,onlytextwidth]
\begin{column}{0.54\textwidth}
Use the Newtonian similarity scaling:
\[
  f_c=\psi\sqrt{\frac{\rho}{\mu_0 xU}},
  \qquad
  \eta = y\sqrt{\frac{\rho U}{\mu_0 x}}.
\]
The remaining streamwise dependence enters through
\[
  \delta = \frac{\lambda^2U^3\rho}{\mu_0x}
  =\frac{\De^2}{\Reyn_{n,x}}.
\]
\end{column}
\begin{column}{0.43\textwidth}
\[
\left(1+\alpha\delta f_c''^2\right)
\left(1+\delta f_c''^2\right)^{(\alpha-3)/2}f_c'''
+\frac{1}{2}f_cf_c''=0.
\]
\begin{itemize}
  \item Solve an ODE at each value of $\delta$.
  \item Large $x$ means small $\delta$: Newtonian behavior.
  \item Near the leading edge, $\delta$ is large: power-law behavior.
\end{itemize}
\end{column}
\end{columns}
\end{frame}

% 10
\begin{frame}{New Sakiadis power-law solutions: what was missing}
\begin{itemize}
  \item Earlier work concluded that Sakiadis power-law solutions fail for $\alpha<1/2$ and $\alpha>1$.
  \item That conclusion depends on assuming a far-field form with $f_p\to C$.
  \item The boundary condition only requires $f_p'\to0$; it does not require $f_p$ to approach a constant.
  \item Alternative far-field structures produce admissible separatrix solutions.
\end{itemize}
\vspace{0.3em}
\begin{center}
\large These new solutions are essential for comparing Sakiadis power-law predictions to Carreau predictions across shear-thinning and shear-thickening regimes.
\end{center}
\end{frame}

% 11
\begin{frame}{Shear-thinning Sakiadis solutions for $0<\alpha<1/2$}
\begin{columns}[T,onlytextwidth]
\begin{column}{0.48\textwidth}
The new far-field balance is
\[
  f_p(\eta;\alpha)\sim H(\eta+J)^p,
  \qquad
  p=\frac{1-2\alpha}{2-\alpha}.
\]
\[
  H=\left(\frac{\alpha(\alpha+1)(2-p)}{p^{1-\alpha}(1-p)^{1-\alpha}}\right)^{1/(2-\alpha)}.
\]
\begin{itemize}
  \item Then $0<p<1$, so $f_p'\to0$ while $f_p$ grows sublinearly.
  \item This satisfies the Sakiadis far-field velocity condition.
\end{itemize}
\end{column}
\begin{column}{0.48\textwidth}
\centering
\figinclude[width=0.94\linewidth]{new_thinning_asymptotic.eps}
\end{column}
\end{columns}
\end{frame}

% 12
\begin{frame}{The admissible solution is a separatrix}
\begin{columns}[T,onlytextwidth]
\begin{column}{0.45\textwidth}
Shoot with
\[
  f_p''(0;\alpha)=\kappa.
\]
There is a critical value $\kappa=\kappa_p$ such that:
\begin{itemize}
  \item above it, solutions remain monotone;
  \item below it, solutions turn over;
  \item at it, $f_p'\to0$ smoothly.
\end{itemize}
\end{column}
\begin{column}{0.52\textwidth}
\centering
\figinclude[width=0.88\linewidth]{fig3.eps}
\end{column}
\end{columns}
\end{frame}

% 13
\begin{frame}{Shear-thickening Sakiadis solutions for $\alpha>1$}
\begin{columns}[T,onlytextwidth]
\begin{column}{0.47\textwidth}
For $\alpha>1$, the separatrix has finite thickness:
\[
  f_p(\eta;\alpha)=C_2 \quad \text{for } \eta\ge C_1.
\]
Consequently,
\[
  u_x=Uf_p'(\eta;\alpha)=0 \quad \text{outside the boundary layer.}
\]
The solution is $C^4$ but not $C^5$ at the boundary-layer edge.
\end{column}
\begin{column}{0.5\textwidth}
\centering
\figinclude[width=0.82\linewidth]{fig6.eps}
\end{column}
\end{columns}
\end{frame}

% 14
\begin{frame}{Finite-width thickening solutions have a fifth-derivative jump}
\begin{columns}[T,onlytextwidth]
\begin{column}{0.48\textwidth}
\begin{itemize}
  \item Velocity, velocity gradient, and curvature remain continuous.
  \item The fifth derivative jumps at the edge of the finite-thickness layer.
  \item This mirrors the known singular structure in shear-thickening Blasius flow.
  \item The Carreau model smooths this because it recovers Newtonian viscosity at low shear.
\end{itemize}
\end{column}
\begin{column}{0.48\textwidth}
\centering
\figinclude[width=0.9\linewidth]{new_thickening_deriv.eps}
\end{column}
\end{columns}
\end{frame}

% 15
\begin{frame}{The new branches connect continuously in wall shear}
\begin{columns}[T,onlytextwidth]
\begin{column}{0.42\textwidth}
Define the dimensionless wall-shear coefficient
\[
  \kappa_p=f_p''(0;\alpha).
\]
The computed values vary continuously through:
\begin{itemize}
  \item the known range $1/2<\alpha\le1$,
  \item the new shear-thinning range $\alpha<1/2$,
  \item the new shear-thickening range $\alpha>1$.
\end{itemize}
\end{column}
\begin{column}{0.55\textwidth}
\centering
\figinclude[width=0.92\linewidth]{fig5.eps}
\end{column}
\end{columns}
\end{frame}

% 16
\begin{frame}{Wall-shear scalings for comparison}
\begin{columns}[T,onlytextwidth]
\begin{column}{0.49\textwidth}
With $\bar\tau_i=\tau_i/(\rho U^2)$,
\[
  \bar\tau_n = \kappa_n \Reyn_{n,x}^{-1/2}.
\]
\[
  \bar\tau_p = \kappa_p^\alpha
  \De^{(\alpha-1)/(\alpha+1)}
  \Reyn_{n,x}^{-\alpha/(\alpha+1)}.
\]
\end{column}
\begin{column}{0.49\textwidth}
\[
  \bar\tau_c = \kappa_c\Reyn_{n,x}^{-1/2}
  \left(1+\kappa_c^2\delta\right)^{(\alpha-1)/2},
  \qquad \kappa_c=f_c''(0;\delta).
\]
\[
  \delta=\frac{\De^2}{\Reyn_{n,x}}.
\]
\begin{block}{Interpretation}
Changing $x$ changes $\Reyn_{n,x}$ and hence changes which part of the rheological curve the boundary layer samples.
\end{block}
\end{column}
\end{columns}
\end{frame}

% 17
\begin{frame}{Carreau wall shear follows the two self-similar asymptotes}
\begin{center}
\figinclude[height=0.36\textheight]{new_shear.eps}\hspace{1em}
\figinclude[height=0.36\textheight]{new_shear_blas.eps}
\end{center}
\vspace{-0.4em}
\begin{itemize}
  \item Near the leading edge: high shear, Carreau follows the power-law similarity prediction.
  \item Far downstream: low shear, Carreau approaches Newtonian wall shear.
  \item The transition occurs where the Newtonian and power-law asymptotes cross.
\end{itemize}
\end{frame}

% 18
\begin{frame}{Asymptotic agreement and a finite high-shear correction}
\begin{columns}[T,onlytextwidth]
\begin{column}{0.5\textwidth}
Downstream:
\[
  \bar\tau_c\sim \bar\tau_n
  \qquad \text{as }\Reyn_{n,x}\to\infty.
\]
Near the leading edge:
\[
  \bar\tau_c\sim (1-c(\alpha))\bar\tau_p
  \qquad \text{as }\Reyn_{n,x}\to0.
\]
\begin{itemize}
  \item The correction is independent of $\De$ in the high-shear limit.
  \item For $\alpha\in[0.3,1.4]$, the high-shear offset is roughly at most $10$--$11\%$.
\end{itemize}
\end{column}
\begin{column}{0.47\textwidth}
\centering
\figinclude[width=0.9\linewidth]{rel_err.eps}\vspace{0.2em}
\figinclude[width=0.9\linewidth]{fig7.eps}
\end{column}
\end{columns}
\end{frame}

% 19
\begin{frame}{Why is there any high-shear error?}
\begin{columns}[T,onlytextwidth]
\begin{column}{0.52\textwidth}
The wall shear is not purely local. For Blasius,
\[
  \bar\tau_i = \frac{\dd}{\dd x}\int_0^\infty u_x(U-u_x)\,\dd y,
  \qquad i=n,c,p.
\]
For Sakiadis, use the analogous expression with $U=0$ in the far field.
\begin{itemize}
  \item The full velocity profile contributes.
  \item Power-law and Carreau profiles can have similar near-wall slopes but different far-field behavior.
  \item That profile difference produces $c(\alpha)$.
\end{itemize}
\end{column}
\begin{column}{0.45\textwidth}
\centering
\figinclude[width=0.98\linewidth]{new_profiles_blas.eps}\vspace{0.2em}
\figinclude[width=0.98\linewidth]{new_profiles.eps}
\end{column}
\end{columns}
\end{frame}

% 20
\begin{frame}{A practical criterion: when to switch models}
Set the Newtonian and power-law wall-shear predictions equal. The crossover occurs at
\[
  \delta^*=
  \left(\frac{\kappa_p^\alpha}{\kappa_n}\right)^{2(1+\alpha)/(1-\alpha)}.
\]
\begin{columns}[T,onlytextwidth]
\begin{column}{0.5\textwidth}
\begin{block}{Rule of thumb}
\begin{itemize}
  \item $\delta/\delta^*\gg1$: use the power-law similarity solution.
  \item $\delta/\delta^*\ll1$: use the Newtonian similarity solution.
  \item Transition region: solve the Carreau one-parameter ODE if accuracy matters.
\end{itemize}
\end{block}
\end{column}
\begin{column}{0.46\textwidth}
\begin{itemize}
  \item $\delta^*$ is relatively insensitive to $\alpha$: roughly $14$--$20$ over a broad range.
  \item Drag estimates can be obtained by splitting the plate at the $x$ corresponding to $\delta=\delta^*$.
\end{itemize}
\end{column}
\end{columns}
\end{frame}

% 21
\begin{frame}{Conclusions}
\begin{enumerate}
  \item Carreau wall shear is well captured by self-similar Newtonian and power-law asymptotes outside the transition region.
  \item The high-shear power-law approximation has a finite, quantified correction $c(\alpha)$ caused by full-profile differences.
  \item New Sakiadis power-law similarity solutions exist for $\alpha<1/2$ and $\alpha>1$.
  \item For $\alpha>1$, the Sakiadis solution is finite-width and $C^4$ but not $C^5$ at its edge.
  \item The crossover parameter $\delta^*$ gives a low-cost practical approximation framework for non-self-similar Carreau boundary layers.
\end{enumerate}
\end{frame}

% 22
\begin{frame}{Backup: minimal numerical recipe}
\begin{itemize}
  \item Convert the relevant BVP to a shooting IVP by prescribing $f''(0)$.
  \item Truncate $\eta\in[0,\infty)$ to $[0,L]$.
  \item Use RK45 integration and choose $f''(0)$ so that the far-field condition is met at $\eta=L$.
  \item Increase $L$ until wall-shear coefficients are converged to the desired tolerance.
\end{itemize}
\vspace{0.4em}
\begin{center}
This is cheap for power-law/Newtonian models and still manageable for Carreau because only the parameter $\delta$ is swept.
\end{center}
\end{frame}

\end{document}

